%! Author = chtompki
%! Date = 1/14/26
% Example LaTeX document for GP111 - note % sign indicates a comment
\documentclass[12pt]{amsart}
% Default margins are too wide all the way around. I reset them here
\setlength{\hoffset}{-.75in}
\setlength{\textwidth}{6.5in}
\setlength{\voffset}{-.5in}
\setlength{\textheight}{9.0in}
\usepackage{amsmath}
\usepackage{amssymb}
\usepackage{amsthm}
\usepackage{pgfplots}
\usepackage{enumerate}
\usepackage{tikz}
\usetikzlibrary{shadows,arrows,arrows.meta,shapes,positioning,calc}% Required for the Circle and Triangle arrow tips
\usepackage{setspace}
\usepackage{siunitx}

\theoremstyle{definition}
\newtheorem{definition}{Definition}

\theoremstyle{question}
\newtheorem{question}{Question}

\title{Seminar Summary:\\March 17, 2026 - Cheng Ly\\Earning a Ph.D.}
\author{Rob Tompkins}
\date{\today}


\begin{document}
    \maketitle
    \date{\today}
    \begin{onehalfspace}
        We began the lecture talking about Neuroscience and talking about how the nevous system is electrical
        in nature.
        Clearly our nevous system has something electrical going on with it, because for example defibrillators
        rely on the application of a voltage potential to help reset the heart.
        We then looked at another example where only electrical stimuli are harnessed in a way as to control a robotic
        apparatus (an arm for example).
        So not only are electrical potentials used for bodily regulation, they are in fact used in all fascets of
        both the somatic and autonomic nervous system (willful control versus involuntary functions).

        After the introduction of the electrical nature of the nervous system, we went into discussing various different
        collaborators of the student variety that Dr. Ly had.


    \end{onehalfspace}
\end{document}
